\documentclass[conference]{IEEEtran}
\IEEEoverridecommandlockouts
\usepackage{cite}
\usepackage{amsmath,amssymb,amsfonts}
\usepackage{algorithmic}
\usepackage{float}
\usepackage{comment}
\usepackage{graphicx}
\usepackage{pdflscape}
\usepackage{siunitx}
\usepackage{textcomp}
\usepackage{rotating}
\usepackage{xcolor}
\def\BibTeX{{\rm B\kern-.05em{\sc i\kern-.025em b}\kern-.08em
    T\kern-.1667em\lower.7ex\hbox{E}\kern-.125emX}}

% =============================================================================
%  GUIA DE FIGURAS - usa APENAS as 4 imagens anexadas. Salve cada uma em imgs/:
%   Fig.1  imgs/field.png            -> campo top-view (TITAN vs VIPER + bola)
%   Fig.2  imgs/objs_scale.png       -> comparacao Viper, Titan e bola
%   Fig.3  imgs/monitor_training.png -> monitor de treino (Viper topo / Titan base):
%                                       recompensa + duracao do episodio por passos
%   Fig.4  imgs/learning_curve.png   -> curva por episodio / eficiencia / taxa de gol
%                                       (PROVA do reward hacking)
% =============================================================================

\begin{document}
\title{Deep Reinforcement Learning for Competitive 1v1 Robot Soccer Using
Heterogeneous Three-Wheel Omni-Directional Robots}

\author{\IEEEauthorblockN{Daniela Osório}
\IEEEauthorblockA{\textit{Faculty of Engineering,} \\
\textit{University of Porto} \\
Porto, Portugal \\
up202208679@up.pt}
\and
\IEEEauthorblockN{Diego Ferreira}
\IEEEauthorblockA{\textit{Faculty of Engineering,} \\
\textit{University of Porto} \\
Porto, Portugal \\
up202205295@up.pt}
\and
\IEEEauthorblockN{Marcella Andrade}
\IEEEauthorblockA{\textit{Faculty of Engineering,} \\
\textit{University of Porto} \\
São Paulo, Brasil \\
up2025121342@up.pt}
\and
\IEEEauthorblockN{Pabline Nascimento}
\IEEEauthorblockA{\textit{Faculty of Engineering,} \\
\textit{University of Porto} \\
Divinópolis, Brasil \\
up202512028@up.pt}
}
\maketitle

% -------- %
% ABSTRACT %
% -------- %
\begin{abstract}
This paper investigates the use of Deep Reinforcement Learning (DRL) for
autonomous robot soccer in a one-versus-one scenario involving heterogeneous
omni-directional mobile robots. Two platforms, Titan and Viper, differing in
size, wheel geometry, and maximum velocity, were trained with Proximal Policy
Optimization (PPO) in a Webots~R2025a environment integrated with Gymnasium and
Stable-Baselines3. A single morphology-conditioned policy controls both robots,
operating on a privileged state representation expressed entirely in the robot's
local reference frame, with the three-wheel omni-directional inverse kinematics
embedded in the low-level controller. Training follows a two-stage Curriculum
Learning (CL) strategy: a foundational single-agent stage (1v0) for locomotion,
ball approach and shooting, followed by an adversarial stage (1v1) that
fine-tunes the policy against an opponent. We report results honestly: the agent
reliably reaches and scores from short and medium range, but a competency-based
analysis exposes a finishing limitation — the deterministic policy converges to a
reward-shaping local optimum in which it parks near a goal post and stalls instead
of completing the shot, yielding a 1v0 goal rate of 37.5\% (with no own-goals or
out-of-bounds) and draws in 1v1. We analyse this behaviour, quantify it, and
identify the reward-balance and training-budget changes required to overcome it.
The study also shows that a robot-local observation design generalises across
platforms with different mass, geometry and speed.
\end{abstract}

\begin{IEEEkeywords}
Robot Soccer, Deep Reinforcement Learning, Proximal Policy Optimization,
Curriculum Learning, Heterogeneous Robots, Omni-Directional Robots,
Autonomous Agents
\end{IEEEkeywords}

% ----- %
% PAPER %
% ----- %
\section{Introduction}
Robot soccer is a well-established benchmark problem in autonomous robotics and
artificial intelligence, as it integrates navigation, manipulation,
decision-making, and adversarial interaction in a dynamic and partially
competitive environment. Learning effective soccer behaviour is particularly
challenging, as it requires agents to jointly acquire locomotion control and ball
manipulation skills while adapting to the presence of an opposing player.

Recent advances in Deep Reinforcement Learning (DRL) have enabled the acquisition
of complex control policies through interaction with simulated environments. In
this work, Proximal Policy Optimization (PPO) and Curriculum Learning (CL) are
employed to train two heterogeneous three-wheel omni-directional robots, Titan and
Viper, in a one-versus-one soccer scenario. The two platforms differ in physical
dimensions and dynamic capabilities, resulting in asymmetric conditions that
necessitate adaptive policy learning. A distinctive aspect of our setup is that a
single PPO policy controls both robots: morphology is encoded as an observation
feature rather than learned by separate models.

The objective of this study is to evaluate the capability of DRL methods to learn
effective behaviours for navigation, ball control, and goal scoring in a
competitive setting, and to characterise the failure modes that arise from dense
reward shaping under a limited training budget. The remainder of this paper is
organised as follows: Section~\ref{sec:related} presents related work,
Section~\ref{sec:method} describes the methodology, Section~\ref{sec:setup}
details the experimental setup, Section~\ref{sec:results} discusses the results,
and Section~\ref{sec:conclusion} concludes the paper and outlines future work.

\section{Related Work} \label{sec:related}
\subsection{Robot Soccer Competitions}
Robot soccer has long served as a benchmark for autonomous robotics because it
combines perception, motion control, decision making, and adversarial interaction
within a single task. Since the introduction of RoboCup, the field has become one
of the most established research platforms for evaluating intelligent robotic
systems~\cite{kitano1998robocup}. RoboCup has expanded into multiple leagues,
including the Small Size League, Humanoid League, and Standard Platform League,
each emphasising different challenges such as high-speed multi-robot coordination,
humanoid locomotion, and software-focused competition under standardised hardware.
In contrast to settings that emphasise homogeneous teams or multi-agent
cooperation, this work focuses on a simplified one-versus-one scenario using
heterogeneous omni-directional robots, allowing the influence of robot morphology
on learned behaviour to be investigated.

\subsection{Reinforcement Learning in Robot Soccer}
Reinforcement Learning (RL) has been widely investigated as a means of enabling
robots to acquire soccer skills directly through interaction with the environment.
Early work such as the RoboCup Keepaway task demonstrated that RL could learn
cooperative behaviours in multi-agent settings~\cite{stone2005keepaway}. More
recently, advances in DRL and simulation frameworks such as rSoccer have enabled
increasingly complex tasks under standardised conditions~\cite{martins2021rsoccer}.
Large-scale studies have shown that DRL can learn agile locomotion, ball handling,
and strategic behaviours for humanoid robots, including policies that transfer to
physical platforms~\cite{haarnoja2023agilesoccer, tirumala2024visionsoccer}. While
these works focus on humanoid locomotion or homogeneous agents, the present work
investigates whether a single PPO policy can adapt to two heterogeneous
omni-directional robots with distinct physical characteristics.

\subsection{Curriculum Learning in Robotics}
Curriculum Learning was introduced as a strategy for improving learning by
presenting tasks in order of increasing difficulty~\cite{bengio2009curriculum}. In
robotic RL, curriculum-based training improves convergence and facilitates the
acquisition of behaviours that are difficult to learn from sparse rewards alone.
Robot soccer naturally benefits from this approach because successful play requires
the sequential acquisition of navigation, ball manipulation, and opponent
interaction. Following this principle, our methodology separates training into a
foundational 1v0 stage and a competitive 1v1 stage.

\section{Methodology} \label{sec:method}
\subsection{Simulation Environment}
All experiments were conducted in Webots~R2025a, an open-source robot simulator
that provides physics-based rigid body dynamics, sensor simulation, and a Python
API. The training loop was implemented using the Gymnasium interface, with
Stable-Baselines3 providing the PPO implementation. The supervisor controller is a
privileged Webots process with omniscient access to all node positions; it serves
as the Gymnasium environment, driving the simulation step by step and exposing the
RL interface.

The field was modelled as a scaled-down rectangular pitch (scaling factor 0.65
relative to standard dimensions) measuring $7.4 \times 10.4\,\mathrm{m}$, using the
Webots NUE coordinate convention (X east, Y up, Z south). Two 1.5\,m-wide goals
were positioned at opposite ends along the Z-axis ($z = \pm 4.55\,\mathrm{m}$). A
spherical ball (radius 0.025\,m, mass 0.055\,kg) with realistic friction and
restitution was placed as a dynamic rigid body. The complete setup is shown in
Fig.~\ref{fig:field_scale}. Each simulation step corresponds to 5 Webots basic
timesteps of 8\,ms, yielding a 40\,ms control period (25\,Hz). Training ran in
\texttt{SIMULATION\_MODE\_FAST}, decoupled from real time, to maximise throughput.
An episode terminates when (i)~the ball enters a goal, (ii)~the ball leaves the
field, or (iii)~the maximum number of simulation steps is reached.

\begin{table}[h]
    \centering
    \caption{Simulation environment parameters.}
    \label{tab:environment}
    \begin{tabular}{lc}
        \hline
        \textbf{Parameter} & \textbf{Value} \\
        \hline
        Field dimensions   & $7.4 \times 10.4$ m \\
        Coordinate system  & Webots NUE \\
        Goal width         & 1.5 m \\
        Ball radius        & 0.025 m \\
        Control period     & 40 ms (25 Hz) \\
        Physics timestep   & 8 ms \\
        \hline
    \end{tabular}
\end{table}

% FIG.1 - imagem anexada: campo top-view
\begin{figure}[H]
    \centering
    \includegraphics[width=\linewidth]{imgs/field.png}
    \caption{Omnidirectional robot soccer 1v1 field: Viper (red) and Titan (blue)
    with the ball at the centre; goals at opposite ends along the Z-axis.}
    \label{fig:field_scale}
\end{figure}

\subsection{Robot Models}
Two holonomic platforms were used: Viper and Titan, both implemented as Webots
PROTO nodes with omnidirectional drive. They share the same sensor suite
(1440-ray LiDAR, IMU, GPS) but differ in physical parameters, as illustrated in
Fig.~\ref{fig:robot_scale} and summarised in Table~\ref{tab:robots}. Viper is
smaller and faster; Titan is larger and more robust.

\begin{table}[h]
    \centering
    \caption{Physical and sensor parameters of the two robot platforms.}
    \label{tab:robots}
    \begin{tabular}{lcc}
        \hline
        \textbf{Parameter} & \textbf{Viper} & \textbf{Titan} \\
        \hline
        Colour & \textcolor{red}{Red} & \textcolor{blue}{Blue} \\
        Wheel radius (m) & 0.030 & 0.035 \\
        Wheel distance (m) & 0.138 & 0.170 \\
        Max wheel velocity (rad/s) & 15.0 & 9.0 \\
        Approx. top speed (m/s) & 0.45 & 0.315 \\
        Body radius (m) & 0.120 & 0.165 \\
        LiDAR rays & 1440 & 1440 \\
        LiDAR max range (m) & 3.50 & 3.50 \\
        \hline
    \end{tabular}
\end{table}

% FIG.2 - imagem anexada: comparacao de tamanho
\begin{figure}[H]
    \centering
    \includegraphics[width=\linewidth]{imgs/objs_scale.png}
    \caption{Viper (red), Titan (blue), and ball size comparison.}
    \label{fig:robot_scale}
\end{figure}

\subsection{Three-Wheel Omnidirectional Inverse Kinematics}
Each robot is driven by three omnidirectional wheels placed at angular offsets
$\varphi_1 = 0^\circ$, $\varphi_2 = 120^\circ$, $\varphi_3 = 240^\circ$ relative to
the forward axis ($+x$ in body frame). Let $v_x$ and $v_z$ denote forward and
lateral translational velocities (m/s) and $\omega$ the yaw rate (rad/s). The
inverse-kinematics mapping from body velocities to wheel angular speeds
$[\omega_1,\omega_2,\omega_3]^\top$ is
\begin{equation}
\begin{bmatrix} \omega_1 \\ \omega_2 \\ \omega_3 \end{bmatrix}
= \frac{1}{R}
\begin{bmatrix}
1 & 0 & L \\
-\sin 30^\circ & \cos 30^\circ & L \\
-\sin 30^\circ & -\cos 30^\circ & L
\end{bmatrix}
\begin{bmatrix} v_x \\ v_z \\ \omega \end{bmatrix},
\end{equation}
where $R$ is the wheel radius and $L$ the centre-to-wheel distance. Embedding the
inverse kinematics in the low-level controller decouples the policy from
wheel-level actuation, so the network outputs intuitive Cartesian body velocities
and each wheel speed is clamped to the robot-specific maximum before application.

\subsection{Observation Space}
The policy operates on a privileged state representation provided by the supervisor,
without direct access to raw sensor measurements. During the 1v0 stage the
observation vector contains 24 dimensions; in the 1v1 stage six additional
opponent-related features are included, yielding 30 dimensions. The complete space
is summarised in Table~\ref{tab:obs}; features 24--29 are present only in 1v1.

The robot yaw $\theta$ is extracted from the orientation matrix as
$\theta = \operatorname{atan2}(-m_1, m_0)$, where $m_0$ and $m_1$ are elements of
the matrix. World-frame direction vectors are transformed into the robot's local
frame by
\begin{equation}
\mathbf{u}_l =
\begin{bmatrix}
\cos\theta & \sin\theta \\
-\sin\theta & \cos\theta
\end{bmatrix}
\mathbf{u}_w ,
\end{equation}
where $\mathbf{u}_w = (d_x, d_z)$ is a world-frame unit direction and
$\mathbf{u}_l$ its local-frame representation. Consequently positive local-frame
$z$-values indicate that the target lies in front of the robot, so the policy maps
$d_z^l > 0 \Rightarrow v_z > 0$ without learning an implicit coordinate rotation.

\begin{table}[t]
    \centering
    \caption{Observation features. Features 24--29 are omitted in 1v0; distances
    are normalised and directions are robot-relative.}
    \label{tab:obs}
    \begin{tabular}{c l c}
        \hline
        \textbf{Index} & \textbf{Feature} & \textbf{Range} \\
        \hline
        0      & Robot type ID                                  & $\{0,1\}$ \\
        1      & Distance to ball                               & $[0,1]$ \\
        2--3   & Direction to ball $(x,z)$                      & $[-1,1]^2$ \\
        4      & Distance to attack goal centre                 & $[0,1]$ \\
        5--6   & Direction to attack goal centre $(x,z)$        & $[-1,1]^2$ \\
        7--9   & Attack right post: distance, direction         & $[0,1]\times[-1,1]^2$ \\
        10--12 & Attack left post: distance, direction          & $[0,1]\times[-1,1]^2$ \\
        13--15 & Own right post: distance, direction            & $[0,1]\times[-1,1]^2$ \\
        16--18 & Own left post: distance, direction             & $[0,1]\times[-1,1]^2$ \\
        19--21 & Robot velocity $(v_x,v_z,\omega)$              & $\mathbb{R}^3$ \\
        22--23 & Ball velocity $(v_x,v_z)$                      & $\mathbb{R}^2$ \\
        24     & Distance to opponent                           & $[0,1]$ \\
        25--26 & Direction to opponent $(x,z)$                  & $[-1,1]^2$ \\
        27--28 & Opponent velocity $(v_x,v_z)$                  & $\mathbb{R}^2$ \\
        29     & Opponent blocking flag                         & $\{0,1\}$ \\
        \hline
    \end{tabular}
\end{table}

When no opponent is present, the opponent features take default values
corresponding to maximum normalised distance and null direction/velocity.

\subsection{Action Space}
The policy outputs a continuous action $\mathbf{a} = [a_0, a_1, a_2] \in [-1,1]^3$
mapped to velocity commands by
\begin{equation}
v_x = a_0\,v_{\max}, \quad v_z = a_1\,v_{\max}, \quad \omega = a_2\,\omega_{\max},
\end{equation}
with $v_{\max} = 0.5\,\mathrm{m/s}$ and $\omega_{\max} = 3.0\,\mathrm{rad/s}$. The
commands are serialised, Base64-encoded, and transmitted to the robot controller
over an IPC channel — Base64 ensures compatibility with the Webots~R2025a receiver
interface, which decodes messages as UTF-8 strings. Producing all three components
simultaneously enables full omnidirectional motion.

\subsection{Reward Design}
Two reward functions were designed: a simpler one for 1v0 and an extended version
for 1v1. Both share the same dense shaping hierarchy and differ only in terminal
signals and opponent-aware terms. The core shaping components are:
\begin{itemize}
    \item \textbf{Per-step time penalty:} $-0.003$.
    \item \textbf{Robot-to-ball progress:} $+4.0 \cdot \Delta d_{\text{robot}\to\text{ball}}$.
    \item \textbf{Ball-to-goal progress:} $+15.0 \cdot \max(0,\Delta d_{\text{ball}\to\text{goal}})$, clipped on backward motion to prevent ball avoidance.
    \item \textbf{Ball velocity toward goal:} up to $\pm1.0$.
    \item \textbf{Robot--ball--goal alignment:} $+0.3 \cdot \cos\theta$.
    \item \textbf{Contact bonus:} $+0.25$ for directed touches ($+0.05$ otherwise), at most once every eight steps.
    \item \textbf{Anti-stall penalties:} $-0.10$/step after 20 motionless steps; $-0.30$/step after 20 steps stuck against a post.
\end{itemize}
Terminal signals: in 1v0, $+300$ for a goal, $-200$ for an own goal, $-25$ for
ball out of bounds. In 1v1, $+300$ for a goal, $-250$ when the opponent scores, and
an opponent-proximity penalty of $-0.08$ when the robots are within $0.30\,\mathrm{m}$.

\subsection{Curriculum Learning Strategy}
Both stages employ a competency-based curriculum that adjusts spawn positions to
progressively increase difficulty, avoiding the sparse-reward problem of fully
random initialisation. In the \textbf{1v0} stage the curriculum defines eight
phases (0--7): early phases place the ball within $0.20\,\mathrm{m}$ of the attack
goal with the robot aligned behind it, and each subsequent phase increases the
ball-to-goal distance until phase~7 spawns both entities uniformly across the
field. The per-phase episode budget grows from 350 to 1250 steps.

An initial promotion criterion (goal rate $\geq 50\%$ over 40 episodes) proved
effectively unreachable beyond the trivial phases and stalled the curriculum; it
was revised to $\geq 40\%$ over 30 episodes, complemented by a demotion rule (drop
one phase below a $10\%$ goal rate) so the agent can re-consolidate. With this
revision the curriculum reliably advances through the mid-distance phases.

In the \textbf{1v1} stage the curriculum has four named phases
(\textit{warmup}/\textit{phase1}/\textit{phase2}/\textit{phase3}). During warmup and
phase1 the opponent is static; phase2 introduces a scripted opponent (go-to-ball,
shoot at its own goal); phase3 uses a frozen snapshot of the opponent's policy,
refreshed each epoch, with fully random spawns. The 1v1 policy is obtained from the
1v0 policy by weight surgery: the input layer is extended from 24 to 30 dimensions,
copying the existing weights and zero-initialising the six new opponent inputs.

\subsection{PPO Training Setup} \label{sec:ppo_setup}
All agents were trained using PPO~\cite{schulman2017ppo} from Stable-Baselines3. A
shared MLP policy ($[256,256]$) was used for both platforms; rather than training
separate networks, a single policy conditions on the \texttt{type\_id} feature.
The final configuration uses $n_{\text{steps}} = 2048$, batch size $256$,
$n_{\text{epochs}} = 10$, $\gamma = 0.99$, GAE $\lambda = 0.95$, clip range $0.2$,
entropy coefficient $0.03$, and a linear learning-rate schedule from
$3\times10^{-4}$ to $1\times10^{-5}$. Rewards are normalised online with
\texttt{VecNormalize} (clip value 50, large enough to preserve the $+300$ goal
signal). In the 1v0 stage, a single policy controls either robot via a hot-swap
mechanism that replaces the physical node in the scene while keeping the DEF name
constant; in the 1v1 stage, the active agent alternates between Viper and Titan
across epochs, with the other robot acting as opponent.

\section{Experimental Setup} \label{sec:setup}
\subsection{Training Configuration}
Table~\ref{tab:training_config} summarises the configuration of each stage.

\begin{table}[h]
    \centering
    \caption{Training configuration per stage.}
    \label{tab:training_config}
    \begin{tabular}{lcc}
        \hline
        \textbf{Parameter} & \textbf{1v0} & \textbf{1v1} \\
        \hline
        Observation dims       & 24             & 30 \\
        Curriculum phases      & 8 (0--7)       & 4 (warmup--phase3) \\
        Promotion threshold    & 40\% / 30 eps  & 60\% / 10 eps \\
        Demotion threshold     & 10\% / 30 eps  & --- \\
        Max steps/episode      & 350--1250      & 1000 \\
        Simulation step        & \SI{40}{ms}    & \SI{40}{ms} \\
        Robot types trained    & Viper + Titan  & Viper + Titan (alternating) \\
        Opponent               & None           & Static $\Rightarrow$ scripted $\Rightarrow$ frozen \\
        \hline
    \end{tabular}
\end{table}

\subsection{Evaluation Protocol}
Trained policies were assessed through repeated episodes under \emph{fixed}
curriculum conditions: the curriculum was frozen during evaluation so the
difficulty could not drift, and deterministic action selection was used. For the
1v0 stage we additionally evaluate per phase (sweeping phases 0--7) and against a
random-action agent as a baseline. The 1v1 stage is evaluated as Viper-versus-Titan
matches reporting wins, losses, and draws. Per-episode metrics are aggregated as
mean and standard deviation.

\subsection{Performance Metrics}
We collect: episode reward; \textbf{goal rate}; \textbf{own-goal} and
\textbf{ball-out} rates; \textbf{timeout rate} (episodes that exhaust the step
budget without a terminal event); time-to-score; and final ball-to-goal distance.
For 1v1 we additionally track opponent goal rate and the win/draw outcome.

\section{Results and Discussion} \label{sec:results}
We report results as measured within the available training budget, including a
degenerate behaviour that is presented as a finding because it directly motivates
the reward-design discussion.

\subsection{Learning Performance}
Fig.~\ref{fig:training_curves} shows the training monitor for both robots
(accumulated reward and episode duration over training steps), and
Fig.~\ref{fig:learning_curve} the per-episode learning curve for Viper. Early
phases confirm the curriculum's effectiveness: with the ball placed a few
centimetres from the goal, the agent scores immediately and is exposed to the
$+300$ terminal signal within the first episodes, giving PPO a strong initial
gradient. As the curriculum advances to harder phases the mean reward remains high
and stable (Viper $\approx$208, Titan $\approx$184; best Viper episode $565$), yet
the episode duration stays long ($\approx$729 steps for Viper and $\approx$767 for
Titan) — most episodes reach the step budget without scoring. Correspondingly, the
per-episode goal rate in Fig.~\ref{fig:learning_curve} \emph{falls} from roughly
$80\%$ in the trivial phases to $20$--$30\%$ and does not recover. Critically,
reward stays high while goals decline: this gap is the quantitative signature of
reward shaping being optimised without the terminal objective being achieved, and
it appears for both embodiments — confirming that the bottleneck is the reward
balance rather than a specific robot.

% FIG.3 - imagem anexada: monitor de treino (Viper topo / Titan base)
\begin{figure}[H]
    \centering
    \includegraphics[width=\linewidth]{imgs/monitor_training.png}
    \caption{Training monitor for Viper (top) and Titan (bottom): accumulated
    episode reward (shaping) and episode duration over training steps. Reward
    stabilises (Viper $\approx$208, Titan $\approx$184) while episodes stay long
    ($\approx$729 and $\approx$767 steps) and goals remain rare — the agents
    optimise the dense shaping without reliably finishing.}
    \label{fig:training_curves}
\end{figure}

% FIG.4 - imagem anexada: curva por episodio (PROVA do reward hacking)
\begin{figure}[H]
    \centering
    \includegraphics[width=\linewidth]{imgs/learning_curve.png}
    \caption{Viper per-episode results: reward and moving average (left), sample
    efficiency over total timesteps (centre), and goal rate (right). The reward
    plateaus while the goal rate decays — evidence of a reward-shaping local
    optimum.}
    \label{fig:learning_curve}
\end{figure}

\subsection{Goal Scoring Performance}
Table~\ref{tab:eval1v0} reports a deterministic 1v0 evaluation with the curriculum
frozen. The agent never concedes own-goals and never sends the ball out, but a
majority of episodes end in timeout: the ball is consistently brought close to the
goal (mean final ball-to-goal distance $0.283 \pm 0.120\,\mathrm{m}$) yet not
driven across the line. The high reward standard deviation ($134.05 \pm 139.46$)
reflects a bimodal outcome — episodes either score (best $+320.77$ in 55 steps) or
time out — rather than gradual performance.

\begin{table}[h]
    \centering
    \caption{1v0 deterministic evaluation (fixed curriculum phase).}
    \label{tab:eval1v0}
    \begin{tabular}{lc}
        \hline
        \textbf{Metric} & \textbf{Value} \\
        \hline
        Goal rate            & $37.5\%$ (3/8) \\
        Timeout rate         & $62.5\%$ (5/8) \\
        Own-goal rate        & $0.0\%$ \\
        Ball-out rate        & $0.0\%$ \\
        Mean episode reward  & $134.05 \pm 139.46$ \\
        Mean episode length  & $189.6 \pm 93.4$ steps \\
        Mean final ball-to-goal dist. & $0.283 \pm 0.120$ m \\
        Fastest goal         & $55$ steps ($R = +320.77$) \\
        \hline
    \end{tabular}
\end{table}

\subsection{Ball Control Behaviour and the Park-Near-Post Failure Mode}
Qualitative inspection of the deterministic rollouts reveals the cause of the
timeouts: the policy frequently navigates to a fixed point near a goal post and
stalls there instead of completing the shot. This is a local
optimum of the dense reward — the post-relative observation features and the
alignment term make a near-goal static position locally acceptable, while the
anti-stall penalties ($-0.10$/step motionless, $-0.30$/step post-stuck) are too
small to overcome it and the time penalty ($-0.003$) is negligible. The $0\%$
ball-out and own-goal rates indicate a conservative policy that does not lose the
ball but is under-aggressive at finishing.

\subsection{Titan vs Viper Analysis}
Both morphologies were trained with the same shared policy conditioned on
\texttt{type\_id}, confirming that a robot-local observation design generalises
across platforms with different mass, geometry and speed: Titan, despite its lower
top speed and greater mass, follows the same qualitative learning trajectory as
Viper. The faster Viper reaches scoring positions in fewer steps, whereas Titan
requires longer trajectories to position behind the ball, which under a fixed step
budget reduces its scoring opportunities. Both platforms exhibit the same finishing
limitation.

\subsection{1v1 Adversarial Performance}
The 1v1 policy was initialised from the 1v0 weights and fine-tuned through the
four-phase pipeline. Over 20 evaluation episodes in the self-play phase, all
matches ended in draws after the full 1000-step horizon, with neither robot
scoring. This is consistent with the 1v0 finishing weakness compounded by the
opponent obstructing the goal mouth: transferring an under-finishing base policy
into a harder adversarial task produced mutual stalling rather than scoring.

\subsection{Key Design Decisions}
Three decisions were critical to reaching the behaviour reported above.
\emph{(i)~No CNN on LiDAR:} an earlier version processed 360 LiDAR rays with a 1D
CNN and never scored after 611k steps because the contact bonus dominated and the
agent learned to hover; the current policy uses only analytical GPS-derived
features. \emph{(ii)~$\max(0,\Delta)$ progress clamp:} without it, pushing the ball
the wrong way is penalised and the rational policy avoids the ball; with it, any
contact is neutral or positive. \emph{(iii)~Goal-mouth clamped target:} the
ball-to-goal distance uses the nearest point inside the goal mouth, preventing the
degenerate strategy of always pushing toward $x=0$.

\section{Conclusions and Future Work} \label{sec:conclusion}
We presented a two-stage DRL pipeline for 1v1 robot soccer with heterogeneous
omni-directional robots in Webots, using PPO, a shared morphology-conditioned
policy, embedded inverse kinematics, a local-frame observation design, and a
competency-based curriculum. The pipeline produces locomotion and short/medium-range
scoring in 1v0 and a complete 1v1 framework with weight surgery and scripted /
frozen-policy opponents. The main negative finding is a finishing limitation: the
policy brings the ball near the goal and then stalls close to a post, a local
optimum of the dense reward that the current anti-stall penalties do not overcome.
Consequently the honest fixed-phase 1v0 goal rate is $37.5\%$ (with $0\%$
own-goals/ball-outs and $62.5\%$ timeouts), and 1v1 matches end in draws. We
attribute this to the balance of the reward shaping and the limited training
budget, both acknowledged project constraints.

Future work will: (i)~re-balance the reward with an explicit anti-park penalty and
a finishing bonus proportional to ball speed across the goal line; (ii)~train
substantially longer with the now-unblocked curriculum; (iii)~adopt true
simultaneous self-play in 1v1; and (iv)~explore dynamic-ball interception and a 2v2
extension with inter-agent communication, as well as a comparison of PPO against
SAC and DDPG on the same environment.

\section*{Acknowledgment}
The authors thank the teaching team of \emph{Topics in Intelligent Robotics}
(M.IA) at the Faculty of Engineering of the University of Porto (FEUP) for guidance
and feedback throughout this project.

\bibliographystyle{IEEEtran}
\bibliography{refs}
\end{document}
